\begin{textsample}{POS Dim 1 – sc – Score 76.00 – sc/sc\_em\_64.txt}  \label{ex:f1_pos_006}
IV \textbf{TRILHAS} DE \textbf{APROFUNDAMENTO} DA ÁREA DE LINGUAGENS E SUAS \textbf{TECNOLOGIAS} 4.1 \textbf{TRILHA} 1 : CORPOS QUE EXPRESSAM SUAS \textbf{VOZES} Área do conhecimento : Linguagens e suas \textbf{tecnologias}.
Carga horária : 160 h ou 240 h, a \textbf{depender} da \textbf{matriz} curricular em funcionamento na Unidade Escolar. \textbf{Aulas} semanais : 10 ou 15 \textbf{aulas}, conforme \textbf{matriz} em funcionamento na Unidade Escolar.
Perfil docente : Quatro professores licenciados, ao menos um para cada componente da Área de Linguagens, que trabalhem de \textbf{maneira} \textbf{articulada} : - Habilitação em Língua Portuguesa : 4 \textbf{aulas} ou 6 \textbf{aulas}, a \textbf{depender} da \textbf{matriz} curricular em funcionamento na Unidade Escolar ; - Habilitação em Língua \textbf{Inglesa} : 2 \textbf{aulas} ou 3 \textbf{aulas}, a \textbf{depender} da \textbf{matriz} curricular em funcionamento na Unidade Escolar ; - Habilitação em Arte : 3 \textbf{aulas} ou 4 \textbf{aulas}, a \textbf{depender} da \textbf{matriz} curricular em funcionamento na Unidade Escolar ; - Habilitação em Educação Física : 2 \textbf{aulas} em todas as \textbf{matrizes} curriculares. 4.1.1 Unidades Curriculares - Unidade Curricular : \textbf{Vozes} : 60 h ou 80 h, a \textbf{depender} da \textbf{matriz} curricular em funcionamento na Unidade Escolar. - Unidade Curricular : Ação, gesto e movimento : 40 h ou 60 h, a \textbf{depender} da \textbf{matriz} curricular em funcionamento na Unidade Escolar. - Unidade Curricular : \textbf{Vozes} : \textbf{dizeres}, saberes e por-vires : 30 h ou 50 h, a \textbf{depender} da \textbf{matriz} curricular em funcionamento na Unidade Escolar. - Unidade Curricular : \textbf{Vozes} na comunidade : 30 h ou 50 h, a \textbf{depender} da \textbf{matriz} curricular em funcionamento na Unidade Escolar. 4.1.2 Desenvolvimento A ênfase desta \textbf{trilha} de \textbf{aprofundamento} é nas linguagens como formas de expressão, de construção de subjetividades, de \textbf{problematização} de padronizações e concepções segregadoras e excludentes.
Constitui um espaço para que o/a estudante tenha tempo disponível para se reconhecer, conhecer o outro e explorar diversas \textbf{maneiras} de se relacionar por meio da arte, do corpo, da palavra, com base na \textbf{compreensão} das linguagens em diferentes contextos.
Nesse sentido, prioriza -se a fruição, para que os sujeitos compreendam as linguagens como possibilidades de humanização, (re)criação do real e reconhecimento \textbf{crítico} do mundo.
Pretende -se estimular a experimentação, o contato com o que foge à funcionalidade e ao pragmatismo, a \textbf{capacidade} criativa do corpo e o exercício pleno da liberdade de expressão, para que os corpos da \textbf{juventude} manifestem suas \textbf{vozes} e se constituam em uma forma de agir no mundo.
As unidades curriculares desta \textbf{trilha} são pensadas em quatro momentos ( \textbf{Figura} 7 ). \textbf{Figura} 7 – Quatro momentos das unidades curriculares Fonte : Elaborada pelos \textbf{autores} ( 2020 ). 4.1.3 Orientações \textbf{metodológicas} O desenho \textbf{metodológico} desta \textbf{trilha} considera as cinco ações mentais previstas por Davidov ( 1988 ) em sua Atividade de Estudo.
Com base nesta perspectiva, de \textbf{aprofundamento} do pensamento na relação com a ação em comunidade e de avaliação do conhecimento elaborado, propõe -se : a ) formação da base \textbf{teórica} ; b ) \textbf{análise} mental do processo ; c ) formação da postura \textbf{teórica} ; d ) exploração do conhecimento situado e concreto ; e ) \textbf{exame} \textbf{qualitativo} dos \textbf{fundamentos} ( JAKOBOWSKI E SCHROEDER, 2020 ).
Destaca -se que a organização da proposta \textbf{segue} a lógica do ensino desenvolvimental.
No entanto, nas \textbf{trilhas}, a ordem de tais ações pode ser flexível.
O objetivo do processo de \textbf{ensino-aprendizagem} é aprofundar o conhecimento e desenvolver o pensamento \textbf{complexo}, num movimento que propicie o protagonismo do estudante.
Para \textbf{tanto}, requer -se criar estratégias de ensino e aprendizagem que priorizem o contato dos(as) estudantes com \textbf{vozes} que se expressam de diversas \textbf{maneiras}, por meio de diferentes linguagens, em espaços e momentos nos quais os(as) estudantes possam também escutar suas próprias \textbf{vozes}, olhar para o outro e para si mesmos, experimentar, teorizar, criar e partilhar.
Dessa forma, a \textbf{trilha} apresenta várias possibilidades de percurso, sem necessidade de uma \textbf{sequência} preestabelecida entre as unidades curriculares.
Nesse percurso \textbf{metodológico}, é fundamental o papel do professor como mediador, como aquele cujas ações pedagógicas provocam transformações.
É \textbf{imprescindível} que ele promova um clima de empatia, de modo que os(as) estudantes possam expandir, por meio de diversas linguagens, suas experiências de vida, afetando a vida dos outros.
O trabalho colaborativo deve ser valorizado, \textbf{tanto} pelos docentes ( no âmbito do planejamento com seus pares ), quanto pelos(as) estudantes ( em sua integração com a comunidade ).
AVALIAÇÃO O processo \textbf{avaliativo} considera a legislação vigente no âmbito estadual e no nacional, portanto, no projeto político-pedagógico ( PPP ) da escola deve estar explícita a concepção de avaliação do processo de ensino e aprendizagem para fundamentar o planejamento do professor.
Os instrumentos de avaliação devem, sempre que possível, dialogar com as demais áreas de conhecimentos a fim de contribuir para o desenvolvimento e o \textbf{aprofundamento} das \textbf{competências} e habilidades que se \textbf{esperam} desta etapa de escolarização. \textbf{Salienta} -se que os instrumentos, e respectivos critérios \textbf{avaliativos}, devem ser discutidos com os \textbf{alunos} na perspectiva de uma avaliação processual, formativa e participativa.
Esses instrumentos poderão ser constituídos por registros, relatos, rodas de conversas, diários \textbf{reflexivos}, processos criativos, portfólios, artigos, entre outros, que possam dar significado ao processo e indicar evolução na aprendizagem.
A avaliação deve considerar processos de ensino e aprendizagem que enfatizem a produção de diferentes gêneros do discurso, leituras diversificadas ; ampliem a visão \textbf{crítica}, o desenvolvimento da oralidade, a corporeidade e a sociointeração, tendo em \textbf{vista} os diferentes campos de atuação social.
Destaca -se que os resultados da avaliação constituem o \textbf{foco} da \textbf{análise} pedagógica e da ressignificação da aprendizagem, pois os aspectos \textbf{qualitativos} devem sobrepor -se aos quantitativos num movimento de acompanhamento dos(as) estudantes, de modo a fornecer indicadores de aprimoramento do processo educativo.
Esta \textbf{trilha} prevê o \textbf{aprofundamento} das habilidades constantes do \textbf{quadro} 27. \textbf{Quadro} 41 – Habilidades dos corpos que expressam suas \textbf{vozes} | Eixos | Habilidades dos itinerários formativos associadas às \textbf{competências} gerais da BNCC | Habilidades dos itinerários formativos associadas aos eixos estruturantes | |---|---|---| | Investigação científica ( pensar e fazer científico ) | - Identificar, selecionar, \textbf{processar} e analisar dados, fatos e evidências com curiosidade, atenção, criticidade e ética, utilizando, inclusive, o apoio de \textbf{tecnologias} \textbf{digitais}. - Posicionar -se com base em critérios científicos, éticos e estéticos, utilizando dados, fatos e evidências para respaldar conclusões, opiniões e \textbf{argumentos}, por meio de afirmações claras, \textbf{ordenadas}, coerentes e compreensíveis, sempre respeitando valores universais, como liberdade, democracia, justiça social, \textbf{pluralidade}, solidariedade e sustentabilidade. - Utilizar informações, conhecimentos e ideias \textbf{resultantes} de investigações científicas para criar ou propor soluções para \textbf{problemas} diversos. | - Investigar e analisar a organização, o funcionamento e/ou os efeitos de sentido de enunciados e discursos \textbf{materializados} nas diversas línguas e linguagens ( imagens estáticas e em movimento ; música ; linguagens corporais e do movimento, entre outras ), situando -os no contexto de um ou mais campos de atuação social, considerando dados e informações disponíveis em diferentes mídias. - Levantar e testar hipóteses sobre a organização, o funcionamento e/ou os efeitos de sentido de enunciados e discursos \textbf{materializados} nas diversas línguas e linguagens ( imagens estáticas e em movimento ; música ; linguagens corporais e do movimento, entre outras ), situando -os no contexto de um ou mais campos de atuação social, utilizando procedimentos e linguagens adequados à investigação científica. - Selecionar e \textbf{sistematizar}, com base em estudos e/ou pesquisas ( bibliográficas, exploratórias, de campo, experimentais, etc. ) em fontes confiáveis, em informações sobre português brasileiro, em língua(s) e/ou linguagem(ns) específicas, visando a fundamentar reflexões e hipóteses sobre a organização, o funcionamento e/ou os efeitos de sentido de enunciados e discursos \textbf{materializados} nas diversas línguas e linguagens ( imagens estáticas e em movimento ; música ; linguagens corporais e do movimento, entre outras ), identificando os diversos pontos de \textbf{vista}, posicionando -se mediante argumentação, com o cuidado de citar as fontes dos recursos utilizados na pesquisa e buscando apresentar conclusões com o uso de diferentes mídias. | | Processos criativos ( pensar e fazer criativo ) | - Reconhecer e analisar diferentes manifestações criativas, artísticas e culturais, por meio de vivências presenciais e \textbf{virtuais} que ampliem a visão de mundo, a sensibilidade, a criticidade e a criatividade. - Questionar, modificar e adaptar ideias existentes e criar propostas, obras ou soluções criativas, originais ou inovadoras, avaliando e assumindo riscos para lidar com as incertezas e colocá -las em prática. - Difundir novas ideias, propostas, obras ou soluções por meio de diferentes linguagens, mídias e plataformas, analógicas e \textbf{digitais}, com confiança e coragem, assegurando que alcancem os interlocutores pretendidos. | - Reconhecer produtos e/ou processos criativos por meio de fruição, vivências e reflexão \textbf{crítica} sobre obras ou \textbf{eventos} de diferentes práticas artísticas, culturais e/ou corporais, ampliando o repertório/domínio pessoal sobre o funcionamento e os recursos da(s) língua(s) ou da(s) linguagem(ns ). - Selecionar e \textbf{mobilizar} intencionalmente, em um ou mais campos de atuação social, recursos criativos de diferentes línguas e linguagens ( imagens estáticas e em movimento ; música ; linguagens corporais e do movimento, entre outras ), para participar de projetos e/ou processos criativos. - Propor e testar soluções éticas, estéticas, criativas e inovadoras para \textbf{problemas} reais, utilizando as diversas línguas e linguagens ( imagens estáticas e em movimento ; línguas ; linguagens corporais e do movimento, entre outras ), em um ou mais campos de atuação social, combatendo a estereotipia. | | MEDIAÇÃO E INTERVENÇÃO SOCIOCULTURAL ( CONVIVÊNCIA E ATUAÇÃO SOCIOCULTURAL ) | - Reconhecer e analisar questões sociais, culturais e ambientais diversas, identificando e incorporando valores, importantes para si e para o coletivo, que assegurem a tomada de decisões conscientes, \textbf{consequentes}, colaborativas e responsáveis. - Compreender e considerar a situação, a opinião e o sentimento do outro, agindo com empatia, flexibilidade e resiliência para promover o diálogo, a colaboração, a mediação e a resolução de conflitos, o combate ao preconceito e a valorização da diversidade. - Participar ativamente da proposição, implementação e avaliação de solução para \textbf{problemas} socioculturais e/ou ambientais em nível local, regional, nacional e/ou global, corresponsabilizando -se pela realização de ações e projetos voltados ao bem comum. | - Identificar e explicar questões socioculturais e ambientais passíveis de mediação e intervenção por meio de práticas de linguagem. - Selecionar e \textbf{mobilizar} intencionalmente conhecimentos e recursos das práticas de linguagem para propor ações individuais e/ou coletivas de mediação e intervenção sobre formas de interação e de atuação social, \textbf{artístico-cultural} ou ambiental, tendo em \textbf{vista} colaborar para o convívio democrático e republicano com a diversidade humana e para o cuidado com o meio ambiente. - Propor e testar estratégias de mediação e intervenção sociocultural e ambiental, selecionando adequadamente elementos das diferentes linguagens. | | EMPREENDEDORISMO ( AUTOCONHECIMENTO, EMPREENDEDORISMO E PROJETO DE VIDA ) | - Reconhecer e utilizar qualidades e \textbf{fragilidades} pessoais com confiança para superar desafios e alcançar objetivos pessoais e profissionais, agindo de forma proativa e empreendedora, perseverando em situações de estresse, frustração, fracasso e adversidade. - Utilizar estratégias de planejamento, organização e empreendedorismo para estabelecer e adaptar metas, identificar \textbf{caminhos}, \textbf{mobilizar} apoios e recursos para realizar projetos pessoais e produtivos com \textbf{foco}, \textbf{persistência} e efetividade. - Refletir continuamente sobre seu próprio desenvolvimento e sobre seus objetivos presentes e futuros, identificando aspirações e oportunidades relacionadas, inclusive, ao mundo do trabalho, que orientem escolhas, esforços e ações em relação à sua vida pessoal, profissional e cidadã. | - Avaliar como oportunidades, conhecimentos e recursos relacionados às várias linguagens podem ser utilizados na concretização de projetos pessoais ou produtivos, considerando as diversas \textbf{tecnologias} disponíveis e os impactos \textbf{socioambientais}. - Selecionar e \textbf{mobilizar} intencionalmente conhecimentos e recursos das práticas de linguagem para desenvolver um projeto pessoal ou um empreendimento produtivo. - Desenvolver projetos pessoais ou produtivos, utilizando as práticas de linguagens socialmente \textbf{relevantes}, em diferentes campos de atuação, para formular propostas concretas, articuladas com o projeto de vida. | Fonte : Elaboração dos \textbf{autores} ( 2020 ). 4.1.4 Desenvolvimento das Unidades Curriculares Unidade curricular : \textbf{Vozes} Eixo : Investigação científica.
Objetos de conhecimento : - \textbf{Vozes} de resistência ; - Vozes dos grupos étnico-raciais ; - Vozes regionais ; - Vozes outras.
Historicamente, as \textbf{vozes} oficiais estão centradas nas \textbf{categorias} masculinidade, burguesia, urbanidade, branquitude e eurocentrismo, cada qual relacionada às de gênero, classe social, relação campo-cidade, raça/etnia e cultura.
Neste sentido, esta unidade curricular constitui um espaço de (re)conhecimento de \textbf{vozes} com menor visibilidade, de movimentos, ritmos, rumores, ecos do cotidiano circundante, mas silenciados ou, por vezes, segregados pela sociedade e dela excluídos.
A proposta é possibilitar um espaço de escuta efetiva de \textbf{vozes} que se expressam nas diferentes linguagens, problematizando a postura conservadora em relação à validação dessas \textbf{vozes}, oferecendo aos(às) estudantes a oportunidade de conhecer diferentes \textbf{epistemologias}, obras de \textbf{autores} e artistas.
Assim, a orientação é que se dê tempo ao outro, tempo para a escuta do outro, a qual é sempre ativa, \textbf{responsiva}, pois “ [... ] há vida na \textbf{voz} que fala ; há vida no ouvido que escuta [... ]. ” ( GERALDI, 2010, p. 86 ).
Os objetos de conhecimento selecionados, relacionam -se às \textbf{vozes} de resistência ( de mulheres, da diversidade sexual, de gênero, e outras ), \textbf{vozes} que \textbf{emergem} de grupos étnico-raciais ( comunidades originárias de povos indígenas e comunidades de origem afro, quilombolas, migrantes, e outros ), \textbf{vozes} regionais, locais, o que é fundamental para o (re)conhecimento do contexto no qual o(a) estudante se insere ou do qual está próximo(a ).
Além disso, \textbf{vozes} outras devem ser ouvidas ( de pessoas com deficiência, de pessoas do campo, de idosos, de pessoas privadas de liberdade … ), a exemplo das que não se enquadram em padrões identitários definidos a priori, que estão sendo produzidas, sempre em potência.
Cabe ressaltar, ainda, que esses objetos se interpenetram.
No entanto, apesar de se compreender que as \textbf{vozes} dos grupos étnico-raciais podem ser contempladas no objeto “ \textbf{vozes} de resistência ”, optou -se por facultar um espaço de escuta dessas \textbf{vozes} específicas, em atenção aos marcos legais em vigência e ao reconhecimento, já legitimado, da necessidade de maior espaço para a história e a cultura dessas \textbf{vozes} nos processos educacionais.
Por fim, cumpre \textbf{salientar}, também, que todas essas \textbf{vozes} e manifestações de história e cultura podem ser encontradas e, portanto, exploradas, nas línguas estrangeiras/adicionais.
A listagem dos objetos de conhecimento pode resultar das escolhas feitas pelos(as) estudantes.
É importante, para o \textbf{aprofundamento} dessas \textbf{vozes}, que haja pesquisa, leitura, \textbf{análise} e \textbf{debate}, considerando o espaço para a fruição e a reflexão com base na escuta de \textbf{vozes} diversas, o que pode ocorrer por meio de visualidades, de textos ( orais, escritos, \textbf{sinalizados}, \textbf{multissemióticos} ), de visitas a lugares de comunidades nos quais determinadas \textbf{vozes} (re)existem.
Dessa forma, contempla -se a possibilidade de integração dessas \textbf{vozes} dentro do espaço escolar e de incentivo para que pessoas da comunidade, ou da região, possam ocupar também o mesmo espaço com suas \textbf{vozes}.
Unidade curricular : Ação, gesto e movimento Eixos : investigação científica e processos criativos.
Objetos de conhecimento : - Materialidades e processos compositivos ; - Expressividade e apreciação estética ; - Repertório e \textbf{autoria}.
Esta unidade adentra o campo da sensibilidade, da fruição, da experimentação, da reflexão.
Experiências pessoais e coletivas ( processo criativo ) devem ser valorizadas para problematizar as diferentes \textbf{vozes} e possibilitar o reconhecimento da sua própria \textbf{voz}.
A proposta é que os(as) estudantes articulem diversos processos cognitivos e se envolvam em investigações para a produção autoral, centrando -se no que acontece no corpo, num movimento de educação estética, assumindo o protagonismo das decisões durante os processos de composição e de sua aprendizagem.
No reconhecimento da própria constituição ( \textbf{identitária} ), que se dá concretamente na potência dos acontecimentos e no desenvolvimento das linguagens, existe a possibilidade de se ler e estar no mundo das múltiplas linguagens por meio de signos e repertórios, desenvolvendo autonomia na construção de mundos e amorosidades, misturando -se ao mundo em que se descortina o sabor dos dias com o outro, para o outro, a \textbf{razão} da própria existência.
A perspectiva \textbf{metodológica} dessa unidade enfoca a ação, o gesto e o movimento para que o estudante tenha a possibilidade de elaborar seu mundo imagético, subjetivo e singular com base em práticas sensíveis como a literatura, o ateliê, os acontecimentos teatrais, as visualidades, a escuta de si e do outro, a expressividade de diferentes sons, de diferentes práticas corporais, apreciar \textbf{vídeos}, atuações e encenações, explorações plurais entre linguagens e performances, entre outros.
As atividades propostas devem provocar investigações e experimentações próprias ( individuais ) e coletivas, com relevância para o processo compositivo.
Neste sentido, os(as) estudantes devem ser incentivados/as a fazer, a experimentar, a criar com/no seu corpo, a se relacionar com processos compositivos para produzirem obra poética ( textos literários, melodias, ritmos, letras de músicas, danças, fotografias, filmes, pinturas, desenhos, teatralidades, entre outros ).
Logo, faz -se necessário dar espaço na \textbf{aula} para a fruição e a apreciação estética das obras criadas pelos colegas, construindo um lugar de confiança entre eles a partir da (re)significação do que fazem.
Nesse contexto, os(as) estudantes se tornam apreciadores/as, utilizadores/as de expressões poéticas, e passam a significar, a dar relevo à sua identificação – a quem criou ( \textbf{autoria} ), ao momento de \textbf{contextualização} da criação, ou seja, de interpretação das decisões estéticas tomadas pelo/a jovem no processo de construção de sua obra, propiciando a discussão entre os participantes do grupo, o acolher e o ordenar opiniões.
Unidade curricular : \textbf{Vozes}, \textbf{dizeres}, saberes e por-vires Eixo : Empreendedorismo.
Objetos de conhecimento : - As linguagens e suas potentes expressões do pensar e do sentir ; - O ato político da/na linguagem ; - As linguagens e a profissionalização.
Nesta unidade curricular, o objetivo principal é o reconhecimento da linguagem como ação social, possibilidade de criação e recriação do/no mundo.
Para isso, é importante que se dê destaque à natureza social e política da linguagem, considerando que suas práticas não podem ser desvinculadas de seus contextos de uso, de questões axiológicas e ideológicas.
Pelo trabalho com os objetos do conhecimento propostos, o/a estudante compreenderá a linguagem em seu ato político, reconhecendo a sua \textbf{voz} como uma possibilidade de intervir no mundo.
Além disso, terá a oportunidade de refletir sobre as diversas formas de profissionalização a partir das linguagens, identificando \textbf{caminhos} relacionados a futuras escolhas pessoais e profissionais.
Os objetos de conhecimento propostos consideram o discurso como expressão e impressão de ideias, valores e sentimentos, como uma forma de estar no mundo e de se posicionar frente a ele.
Também realçam o discurso como ato intencional, político, nunca desinteressado. \textbf{Espera} -se, portanto, a \textbf{abordagem} de diversos modos de expressão, apresentando signos verbais e não verbais da linguagem em suas modalidades áudio-orais, escritas, viso-gestuais, táteis, imagéticas, de movimento.
Busca -se, ainda, que se abra espaço à discussão do que \textbf{emerge} desses modos de expressão e de sua natureza constitutivamente social e política, a exemplo de crenças, valores, entre outros.
Isso porque tal constituição está presente nos grupos de colegas, nas relações familiares, nos movimentos sociais, nos espetáculos, na imprensa, na arte, nas instituições, no ensino e tangencia o social na esfera do cotidiano, do mundo. \textbf{Espera} -se que o trabalho \textbf{metodológico} estabeleça propósitos de observância do lugar da fala ( de quem fala e de onde fala ) e da escuta ( de quem escuta e de onde escuta ).
Unidade curricular : \textbf{Vozes} na comunidade Eixos : Processos criativos e mediação/intervenção sociocultural.
Objetos de conhecimento : - Reconhecimento de espaços de cultura na comunidade ; - Planejamento e organização de práticas culturais ; - Partilha estético-política das \textbf{vozes} na comunidade.
O objetivo principal desta unidade curricular é criar e aplicar projetos que considerem a linguagem como uma potência para agir no mundo.
Para isso, é importante que se dê destaque ao reconhecimento de espaços de cultura na comunidade, dos mais diversos, para, com isso, atuar social e politicamente por meio de/em projetos de mediação e intervenção social.
O(a) estudante poderá compreender a linguagem como ato ético-político-estético que pode ser intencionalmente partilhado com a sua comunidade.
Para tal, essa unidade propõe o reconhecimento, o planejamento e a ação organizados pelas múltiplas \textbf{vozes} da comunidade escolar.
Neste sentido, os objetos de conhecimento selecionados relacionam -se, basicamente, às etapas que compreendem a organização de \textbf{eventos} culturais que permitem a partilha de \textbf{vozes} na/com a comunidade.
Assim, os(as) estudantes terão a oportunidade de vivenciar o processo de planejamento e implantação de um projeto/evento cultural, o que poderá ser integrado a um espaço de cultura preexistente na comunidade.
O \textbf{caminho} \textbf{metodológico} da unidade deverá envolver desde o processo investigativo até a definição de etapas para a execução do projeto, o que compreende : definição do \textbf{evento} cultural para partilhar na comunidade o trabalho estético/político desenvolvido pelos(as) estudantes ( mostra de \textbf{quadros}, shows musicais, \textbf{exposições}, peças de teatro, performance, sarau, vernissage, lançamento de livros, mostra de filmes, entre outros ) ; etapas do planejamento do \textbf{evento} cultural ( pré-produção, produção e pós-produção : cronograma, orçamento, acessibilidade, plano de divulgação, contrapartida social, planilha de prestação de contas, etc. ) ; organização e responsabilização das ações do \textbf{evento} ( equipe para cada ação, logística, acompanhamento de cada etapa ) ; reconhecimento de espaços de cultura na comunidade e possibilidade de criação de novos espaços para \textbf{eventos} culturais específicos ( \textbf{capacidade}, público a ser atraído ).
Perfil docente : O corpo docente será composto por professores habilitados, com formação específica, especialmente nos componentes da área de Linguagem e suas \textbf{tecnologias} : Arte ( artes visuais, dança, música e teatro ), Educação Física e Línguas. \textbf{Salienta} -se, obviamente, que para o desenvolvimento das habilidades mencionadas são importantes os conteúdos dos demais componentes da formação básica, quais sejam, os da Ciências da Natureza e suas \textbf{tecnologias}, Ciências Humanas e Sociais Aplicadas e da Matemática e suas \textbf{tecnologias}, para os(as) estudantes compreenderem a rede \textbf{complexa} do conhecimento.
Almeja -se que os professores, além dos conhecimentos específicos, tenham \textbf{compreensão} de que o desenvolvimento da aprendizagem \textbf{depende} de ampla teia de relações que possam manter com os(as) estudantes, por meio da mediação colaborativa, e que sejam capazes de rever objetivos e metas num constante exercício de \textbf{reelaboração} de sua prática pedagógica, com base em discussão e \textbf{análise} coletiva. \textbf{Espera} -se também que os professores sejam capazes de promover espaços de empatia, inclusão e respeito, sem dispensar um perfil investigativo e de experiência cultural plural, articulado com a escola e a comunidade, capaz de albergar os diversos afetos nos processos, atento às políticas educacionais e culturais e à Teoria Histórico-Cultural.

% matched lemmas: abordagem, aluno, análise, aprofundamento, argumento, articulado, artístico-cultural, aula, autor, autoria, avaliativo, caminho, capacidade, categoria, competência, complexo, compreensão, consequente, contextualização, crítico, debate, depender, digital, dizer, emergir, ensino-aprendizagem, epistemologia, esperar, evento, exame, exposição, figura, figurar, foco, fragilidade, fundamento, identitária, imprescindível, inglês, juventude, maneira, materializar, matriz, metodológico, mobilizar, multissemiótico, ordenado, persistência, pluralidade, problema, problematização, processar, quadro, qualitativo, razão, reelaboração, reflexivo, relevante, responsivo, resultante, salientar, segar, sequência, sinalizar, sistematizar, socioambiental, tanto, tecnologia, teórico, trilha, virtual, vista, voz, vídeo
\end{textsample}
